\documentclass[10pt,twocolumn]{article}
\usepackage[T1]{fontenc}
\usepackage[utf8]{inputenc}
\usepackage{lmodern}
\usepackage[margin=1.8cm,top=2.4cm,bottom=2cm,columnsep=0.6cm]{geometry}
\usepackage{fancyhdr}
\usepackage{titlesec}
\usepackage{booktabs}
\usepackage{amsmath,amssymb,amsfonts}
\usepackage{microtype}
\usepackage{xcolor}
\usepackage{mdframed}
\usepackage{parskip}
\usepackage{enumitem}
\usepackage{hyperref}

\definecolor{rulered}{RGB}{180,30,30}
\definecolor{boxgray}{RGB}{245,245,245}
\definecolor{keywordgray}{RGB}{100,100,100}

\pagestyle{fancy}
\fancyhf{}
\fancyhead[L]{\textsc{\small Claude Mag}}
\fancyhead[C]{\small\textit{Machine Learning Research Digest}}
\fancyhead[R]{\small 2026-07-08}
\fancyfoot[C]{\small\thepage}
\renewcommand{\headrulewidth}{0.8pt}

\titleformat{\section}{\normalsize\bfseries\color{rulered}}{}{0pt}{}%
  [\vspace{-4pt}\color{rulered}\rule{\columnwidth}{0.4pt}]
\titlespacing{\section}{0pt}{8pt}{4pt}

\hypersetup{colorlinks=true,urlcolor=rulered,linkcolor=black}

\begin{document}
\setlength{\parindent}{0pt}
\setlength{\parskip}{4pt}

{\small\color{keywordgray}\textbf{Keywords:}
  data attribution \textbullet{}
  video generation \textbullet{}
  optical flow \textbullet{}
  temporal dynamics}\par\vspace{4pt}

{\LARGE\bfseries\raggedright
  Motion Attribution for Video Generation}\par\vspace{4pt}

{\large\itshape\raggedright
  MOTIVE Identifies Training Clips That Teach Video
  Models \emph{How} Things Move}\par\vspace{6pt}

{\small\textbf{By} Xindi Wu, Despoina Paschalidou, Jun Gao, Antonio Torralba,
  Laura Leal-Taixe, Olga Russakovsky, Sanja Fidler \& Jonathan Lorraine
  (NVIDIA, MIT, Princeton)}\par\vspace{2pt}

{\small\color{keywordgray}arXiv:2601.08828 \textbullet{}
  ICML 2026 Oral}
\par\vspace{2pt}\noindent{\color{rulered}\rule{\columnwidth}{0.6pt}}\vspace{6pt}

Data attribution for video generation has focused on visual appearance while
ignoring the dimension users care about most: motion. MOTIVE (MOTIon attribution
for Video gEneration) introduces motion-weighted influence functions that isolate
temporal dynamics from static appearance, identifying which training clips teach
models how objects and cameras move. Fine-tuning on MOTIVE-selected data achieves
a \textbf{74.1\%} human preference win rate over the base model and raises
VBench dynamic degree from 41.3\% to \textbf{47.6\%}.

\begin{mdframed}[backgroundcolor=boxgray,linecolor=rulered,linewidth=1pt,
    innerleftmargin=6pt,innerrightmargin=6pt,
    innertopmargin=6pt,innerbottommargin=6pt]
\textbf{\small AT A GLANCE}\par\vspace{4pt}
\begin{description}[leftmargin=0pt,labelsep=4pt,itemsep=2pt,parsep=0pt,
    font=\small\bfseries]
  \item[Problem:] Existing video data attribution methods measure influence on
    per-frame appearance, not temporal dynamics; data curation cannot target
    motion quality specifically.
  \item[Why it matters:] Motion realism is the primary dimension separating
    high-quality from poor video generation; without motion-specific attribution,
    curation is blind to it.
  \item[Method:] Motion-weighted optical-flow loss masks focus gradient-based
    influence computation on dynamic spatial-temporal regions.
  \item[Result:] 74.1\% human win rate; VBench dynamic degree 47.6\%
    vs.\ 41.3\% (random) and 43.8\% (whole-video attribution).
  \item[Evidence:] Evaluated on CogVideoX fine-tuning with three curated
    datasets; human studies ($n{=}200$); VBench motion metrics.
\end{description}
\end{mdframed}

\section{The Big Picture}

Video generation has advanced dramatically, but temporal realism---the
physical plausibility of motion---consistently lags behind per-frame
visual quality. The gap is a data problem: it is unclear which training
clips are responsible for teaching a model that objects fall, water flows,
or cameras pan smoothly.

Data attribution, which quantifies each training example's influence on
model behavior, is the natural tool for this problem. But existing attribution
methods compute influence on a generic generation loss, dominated by the
high-frequency spatial structure of individual frames. Motion signals---which
are lower-magnitude and distributed across frame pairs---are drowned out.

MOTIVE addresses this by constructing a motion-specific loss function and
computing influence with respect to it.

\section{Why Existing Approaches Fall Short}

Gradient-based influence functions (TRAK, DataInf, EK-FAC) assign each
training example a score based on its gradient alignment with a test-set
loss. For video, the default loss is reconstruction error summed uniformly
over all pixels and frames.

This aggregation is problematic for motion attribution because:
\begin{itemize}[itemsep=2pt,parsep=0pt]
  \item Static backgrounds contribute most pixels and dominate the loss.
  \item Motion signals are typically localized to a small fraction of
    the spatial extent (moving objects, hands, cameras).
  \item Two clips with identical backgrounds but opposite motion patterns
    have near-identical influence scores under whole-video attribution.
\end{itemize}

Empirically, whole-video attribution selects clips that match the test
sequence's scene composition, not its motion patterns, yielding only
marginal improvements in dynamic degree (43.8\% vs.\ 41.3\% random).

\section{Core Mechanism}

MOTIVE constructs a \textbf{motion-weighted loss mask} from optical flow:
pixels with high inter-frame displacement receive high weight, static
pixels receive near-zero weight. Influence computation then operates on
this masked loss, making gradients sensitive to motion-generating
patterns rather than static appearance.

The motion mask $M(x,t)$ is derived from the optical flow field
$f(x,t) \in \mathbb{R}^2$ between frames $t$ and $t{+}1$:
\[
M(x,t) = \mathrm{softmax}_\tau\!\left(\|f(x,t)\|_2\right)
\]
where the softmax is taken over all spatial positions $x$ and times $t$,
with temperature $\tau$ controlling mask sharpness. This reweights the
loss to be high wherever the scene is moving.

\section{Technical Formulation}

The motion-weighted loss for training clip $z_i$ is:
\[
\mathcal{L}_M(z_i;\theta) = \sum_{t=1}^{T}\sum_{x \in \Omega}
  M(x,t)\cdot\ell\!\left(v_\theta(x,t),\,f_t(x)\right)
\]
where $v_\theta(x,t)$ is the model's prediction for pixel $x$ at frame $t$
and $\ell$ is the per-pixel reconstruction loss.

The influence score of training clip $z_i$ on a motion-quality test loss $\mathcal{L}_M(z_\text{test})$ is:
\[
\mathcal{I}(z_i) = -\nabla_\theta \mathcal{L}_M(z_\text{test})^\top
  H_\theta^{-1}\,\nabla_\theta \mathcal{L}_M(z_i)
\]
where $H_\theta = \frac{1}{n}\sum_{j=1}^{n}\nabla^2_\theta\mathcal{L}_M(z_j)$
is the Hessian. The Hessian-vector product $H_\theta^{-1}v$ is approximated
via 3 iterations of LiSSA (Linear time Stochastic Second-order Algorithm),
making MOTIVE scalable to models with billions of parameters.

\section{Learning or Inference Procedure}

\begin{enumerate}[itemsep=2pt,parsep=0pt]
  \item \textbf{Flow estimation:} compute optical flow $f(x,t)$ for all frames
    of each candidate training clip using a pre-trained flow network (RAFT).
  \item \textbf{Mask construction:} compute $M(x,t)$ via spatiotemporal softmax
    over flow magnitudes.
  \item \textbf{LiSSA precomputation:} compute $v_\text{test} = H_\theta^{-1}\nabla_\theta\mathcal{L}_M(z_\text{test})$
    using 3 LiSSA iterations.
  \item \textbf{Influence scoring:} compute $\mathcal{I}(z_i) = -v_\text{test}^\top\nabla_\theta\mathcal{L}_M(z_i)$
    for each candidate clip.
  \item \textbf{Curation:} select top-$k$ clips by influence score for fine-tuning.
  \item \textbf{Fine-tuning:} standard video diffusion fine-tuning on selected clips.
\end{enumerate}

\section{What the Guarantee Says}

MOTIVE provides no formal approximation guarantee on influence estimation,
but establishes empirically that the LiSSA approximation error is negligible
relative to the rank-ordering of training clips: Spearman rank correlation
between 3-iteration and exact influence (computed on small-scale models
where exact inversion is tractable) is 0.94.

\section{Experimental Findings}

\begin{center}
\small
\begin{tabular}{lcc}
\toprule
\textbf{Curation Method} & \textbf{VBench Dyn.} & \textbf{Human Pref.} \\
\midrule
Random selection   & 41.3\% & -- \\
Whole-video attr.  & 43.8\% & 52.1\% \\
\textbf{MOTIVE}    & \textbf{47.6\%} & \textbf{74.1\%} \\
\bottomrule
\end{tabular}
\end{center}

Motion smoothness (VBench) also improves: MOTIVE achieves 96.8\%
vs.\ 95.2\% for random selection. Physical plausibility (human study
on gravity and collision realism) improves by 18 percentage points.

\section{Ablations and Interpretation}

\textbf{Flow threshold:} Setting $\tau \to 0$ (binary motion mask at median
flow magnitude) reduces VBench dynamic degree to 45.1\%; smooth soft masks
are important.

\textbf{LiSSA steps:} 1 step yields 46.1\% (noisy influence estimates);
3 steps yields 47.6\%; 10 steps yields 47.7\% (negligible gain).

\textbf{Test set choice:} Influence is computed relative to a motion-diverse
test set; using a static-scene test set produces data selection equivalent
to whole-video attribution, confirming the motion mask is the active ingredient.

\textbf{Model agnosticism:} Results replicate on CogVideoX-2B and CogVideoX-5B;
the framework is model-architecture agnostic.

\section*{Reference}

Xindi Wu, Despoina Paschalidou, Jun Gao, Antonio Torralba, Laura Leal-Taixe,
Olga Russakovsky, Sanja Fidler, Jonathan Lorraine.
\textbf{Motion Attribution for Video Generation}.
arXiv:2601.08828, January 2026.
\url{https://arxiv.org/abs/2601.08828}

\end{document}
